\documentclass{article}
\usepackage{amsmath, amssymb}
\setlength{\parindent}{0pt}

\begin{document}

\textbf{CS 2050 --- Notes 09/02}

\bigskip

\textbf{Precedence of quantifiers.} $\forall$ and $\exists$ have higher precedence than all logical operators.

\bigskip

\textbf{Nested quantifiers.} A nested quantifier is when one quantifier is within the scope of another.

\medskip

\textbf{Example.} $\forall x \, \exists y \, (x - 2y = 5)$

\bigskip

Everything within the scope of a quantifier can be thought of as a propositional function.

\medskip

The order of quantifiers is important unless all the quantifiers are the same.
\[
\exists x \, \forall y \, P(x,y) \text{ is not logically equivalent to } \forall y \, \exists x \, P(x,y).
\]

\bigskip

\textbf{Definition.} $P \equiv Q$ means $P \leftrightarrow Q$ is a tautology.

\bigskip

\textbf{Example.} $\forall x \, (W(x) \to F(x))$

\medskip

For all people $x$, if $x$ is a wizard, then $x$ has a funny hat. All wizards have a hat.

\bigskip

\textbf{Example.} $\forall x \, (W(x) \land F(x))$

\medskip

For all people $x$, $x$ is a wizard and has a funny hat. Everyone is a wizard and has a funny hat.

\bigskip

\textbf{Example.} $\exists x \, (W(x) \to F(x))$

\medskip

For some person $x$, if $x$ is a wizard, then $x$ has a funny hat.

\bigskip

\textbf{Example.} There is an adventurer in Faerun who could use a bow and a sword.
\[
\exists x \, (P(x) \land Q(x))
\]

\bigskip

\textbf{Example.} There is an adventurer in Faerun who could use a bow but could not use a sword.
\[
\exists x \, (P(x) \land \lnot Q(x))
\]

\bigskip

\textbf{Example.} Every adventurer in Faerun could use a bow or use a sword.
\[
\forall x \, (P(x) \lor Q(x))
\]

\bigskip

\textbf{Example.} Some student in this class has played Dungeons and Dragons.

\medskip

$P(x) = $ $x$ has played D\&D. \quad Domain $=$ students in the class.
\[
\exists x \, P(x)
\]

\bigskip

\textbf{Example.} $\forall x \, \exists y \, (x < y)$

\medskip

For all $x$, there exists a real number $y$ such that $x < y$.

\bigskip

\textbf{Example.} $\exists x \, \forall y \, (x < y)$

\medskip

There exists $x$ such that for all $y$, $x < y$.

\bigskip

\textbf{Example.} $\forall x \, \forall y \, \big( ((x < 0) \land (y < 0)) \to (xy > 0) \big)$

\medskip

For every $x$ and every $y$, if $x < 0$ and $y < 0$, then $xy > 0$.

\end{document}
